\paragraph*{04.05.06.04 Validierung von Projektergebnissen planen und organisieren}
\kcilabel{para:04.05.06.04_ValidierungVonProjektergebnissen}{04.05.06.04}
\label{para:04.05.06.04_ValidierungVonProjektergebnissen}

% Task
\par Die Validierung der \gls{ricefw}-Objekte sowie des \gls{pb} durch relevante Stakeholder und Endanwender im Rahmen eines \gls{uat} stellte eine zwingende Voraussetzung für die Produktivsetzung dar.

% Action
\par Zur strukturierten Planung und Organisation dieser Validierung erarbeitete ich gemeinsam mit dem Programmtestmanager ein umfassendes \gls{e2e}-Testkonzept. Dieses adressierte die spezifischen fachlichen und technischen Anforderungen des \gls{wsr} und definierte klare Verantwortlichkeiten, Teststufen sowie Abnahmekriterien für die beteiligten Stakeholder.

% Result
\par Trotz der Abgenommen \gls{e2e}-Testkonzepts und der Durchführung von Tests im \gls{fb} wurden nicht alle vorgesehenen Tests konsequent umgesetzt. Dies führte zu erheblichen Herausforderungen bei der Validierung der Projektergebnisse im Rahmen des \gls{uat} (vgl. \autoref{fig:R2HStatusSSTJuni2025}).

% Reflect
\par Die unzureichende Umsetzung der definierten \gls{e2e}-Testaktivitäten im \gls{fb} hatte direkte Auswirkungen auf die Validierbarkeit der Projektergebnisse. Durch enge interne Abstimmung, pragmatische Lösungsansätze und priorisierte Nachtests konnten die identifizierten Defizite jedoch behoben werden (vgl. \autoref{tab:statusuebersicht}). Für zukünftige Vorhaben habe ich daraus abgeleitet, Testprozesse verbindlicher zu verankern und deren Einhaltung frühzeitig und konsequent zu steuern, um die Validierung von Projektergebnissen nachhaltig sicherzustellen.